\documentclass[12pt,a4paper]{article}

\usepackage{amsmath}
\usepackage{amssymb}
\usepackage{array}
\usepackage{booktabs}
\usepackage{caption}
\usepackage{float}
\usepackage{geometry}
\usepackage{graphicx}
\usepackage{hyperref}
\usepackage[utf8]{inputenc}
\usepackage{makecell}
\usepackage{tabularx}
\usepackage[T2A]{fontenc}
\usepackage[english,russian]{babel}

\geometry{margin=2.2cm}
\hypersetup{hidelinks}
\graphicspath{{results/}}
\captionsetup{font=normalsize}
\renewcommand{\arraystretch}{1.15}

\newcolumntype{Y}{>{\centering\arraybackslash}X}

\title{Отчет по заданию 4\\Построение оптимального портфеля}
\author{}
\date{}

\begin{document}
\maketitle

\section{Постановка задачи}

В задании нужно построить оптимальный портфель для шести акций. Используется модель Тобина--Шарпа--Линтнера:
\[
GAZP,\ ROSN,\ LKOH,\ FEES,\ SBER,\ VTBR.
\]

Для этих бумаг были оценены ожидаемые доходности и ковариационная матрица. Затем были построены эффективные границы и оптимальные портфели при разных значениях коэффициента неприятия риска $\theta$.

Безрисковый актив учтен в двух вариантах:
\begin{itemize}
    \item вложение свободных средств по ставке ЦБ РФ;
    \item заимствование по ставке ЦБ РФ плюс 5 процентных пунктов.
\end{itemize}

\section{Данные}

Исходный файл с данными: \texttt{Data\_zad4\_2026.xlsx}. Сам Excel-файл не изменялся. Для расчетов данные были аккуратно экспортированы в CSV-файлы в папке \texttt{prepared}.

В тексте задания указан интервал 01.04.2016--10.04.2016, но в выданном Excel-файле фактически лежат даты 01.09.2010--01.10.2010. Поэтому расчет выполнен по реальным данным из файла: 23 строки цен и 22 строки дневных доходностей.

Ставка ЦБ РФ для этого периода принята равной 7.75\% годовых. Для варианта с заимствованием использована ставка:
\[
r_b = 7.75\% + 5\% = 12.75\%.
\]

\begin{figure}[H]
    \centering
    \includegraphics[width=\textwidth]{infographic_market_overview.png}
    \caption{Общий вид данных: динамика цен, накопленная доходность и дневные изменения}
\end{figure}

По графику видно, что интервал короткий, поэтому оценки доходности получаются чувствительными к движению цен внутри этих нескольких недель. Это особенно важно при чтении итоговых процентов доходности.

\section{Оценка доходности и риска}

Для каждой акции были рассчитаны простые дневные доходности:
\[
r_{t,i} = \frac{P_{t,i}}{P_{t-1,i}} - 1.
\]

Дальше дневные оценки были переведены в годовые:
\[
m_i = 252\bar r_i, \qquad C = 252\operatorname{cov}(r).
\]

\begin{table}[H]
    \centering
    \caption{Оценки ожидаемой доходности и риска}
    \begin{tabular}{lcc}
        \toprule
        Бумага & Доходность, годовых & Риск, годовых \\
        \midrule
        GAZP & -8.08\% & 12.98\% \\
        ROSN & 61.95\% & 22.29\% \\
        LKOH & 64.92\% & 22.26\% \\
        FEES & 45.98\% & 17.90\% \\
        SBER & 120.37\% & 24.40\% \\
        VTBR & 88.53\% & 24.74\% \\
        \bottomrule
    \end{tabular}
\end{table}

Наибольшую оценочную доходность показал SBER. У GAZP на данном интервале оценка доходности отрицательная. Поэтому оптимизатор почти всегда старается держать GAZP на минимальной допустимой доле.

\begin{table}[H]
    \centering
    \small
    \caption{Годовая ковариационная матрица}
    \resizebox{\textwidth}{!}{
    \begin{tabular}{lrrrrrr}
        \toprule
        & GAZP & ROSN & LKOH & FEES & SBER & VTBR \\
        \midrule
        GAZP & 0.0169 & 0.0162 & 0.0056 & 0.0146 & 0.0208 & 0.0207 \\
        ROSN & 0.0162 & 0.0497 & 0.0110 & 0.0109 & 0.0361 & 0.0232 \\
        LKOH & 0.0056 & 0.0110 & 0.0496 & 0.0064 & 0.0168 & 0.0223 \\
        FEES & 0.0146 & 0.0109 & 0.0064 & 0.0320 & 0.0177 & 0.0176 \\
        SBER & 0.0208 & 0.0361 & 0.0168 & 0.0177 & 0.0595 & 0.0437 \\
        VTBR & 0.0207 & 0.0232 & 0.0223 & 0.0176 & 0.0437 & 0.0612 \\
        \bottomrule
    \end{tabular}}
\end{table}

\begin{figure}[H]
    \centering
    \includegraphics[width=\textwidth]{infographic_risk_correlation.png}
    \caption{Сравнение риска, доходности и взаимосвязи бумаг}
\end{figure}

Риск у бумаг находится примерно в одном диапазоне, но доходности различаются сильно. Поэтому итоговый портфель в первую очередь ограничивается не риском отдельных бумаг, а заданными ограничениями на доли и группы.

\section{Ограничения}

На каждую бумагу наложено простое ограничение:
\[
0.05 \le x_i \le 0.39.
\]

Также использованы групповые ограничения. Группы заданы так:
\[
\begin{array}{ll}
\text{нефтегаз:} & GAZP,\ ROSN,\ LKOH,\\
\text{энергетика:} & FEES,\\
\text{банки:} & SBER,\ VTBR,\\
\text{внутренний рынок:} & LKOH,\ FEES,\ SBER,\ VTBR,\\
\text{внешний рынок:} & GAZP,\ ROSN.
\end{array}
\]

\begin{table}[H]
    \centering
    \caption{Групповые ограничения}
    \begin{tabular}{lcc}
        \toprule
        Группа & Нижняя граница & Верхняя граница \\
        \midrule
        Нефтегаз & 25\% & 65\% \\
        Энергетика & 27\% & 75\% \\
        Банки & 15\% & 55\% \\
        Внутренний рынок & 25\% & 85\% \\
        Внешний рынок & 10\% & 35\% \\
        \bottomrule
    \end{tabular}
\end{table}

Для третьего типа ограничений, \texttt{GroupComparison}, в тексте задания даны формулы, но не указаны конкретные числа. Поэтому в расчете использован простой вариант:
\begin{itemize}
    \item внутренний рынок / внешний рынок: от 1.0 до 5.0;
    \item нефтегаз / банки: от 0.8 до 3.0.
\end{itemize}

Эти значения можно заменить в MATLAB-скрипте, если нужны другие границы.

\section{Модель}

Пусть $x$ --- вектор долей в рискованных активах. Тогда доля безрискового актива равна:
\[
x_0 = 1 - \sum_i x_i.
\]

Если $x_0 > 0$, часть капитала лежит в безрисковом активе. Если $x_0 < 0$, портфель использует заемные средства. Для заданного $\theta$ максимизировалась функция:
\[
U(x) = r_f x_0 + m^\top x - \frac{\theta}{2}x^\top Cx.
\]

Смысл простой: инвестор хочет больше доходности, но платит штраф за риск. Чем больше $\theta$, тем осторожнее инвестор.

\section{Результаты}

\subsection{Сравнение наборов ограничений}

Сначала посмотрим, как меняется портфель при $\theta = 3$.

\begin{table}[H]
    \centering
    \small
    \caption{Влияние ограничений на результат при $\theta = 3$}
    \resizebox{\textwidth}{!}{
    \begin{tabular}{lccccccccc}
        \toprule
        Ограничения & Доходность & Риск & $\xi$ & GAZP & ROSN & LKOH & FEES & SBER & VTBR \\
        \midrule
        Только индивидуальные & 135.73\% & 32.25\% & 2.00 & 5.00\% & 39.00\% & 39.00\% & 39.00\% & 39.00\% & 39.00\% \\
        Индивидуальные + групповые & 90.63\% & 20.67\% & 1.20 & 5.00\% & 30.00\% & 5.00\% & 27.00\% & 39.00\% & 14.00\% \\
        Все ограничения & 90.32\% & 20.53\% & 1.20 & 5.00\% & 30.00\% & 6.33\% & 27.00\% & 39.00\% & 12.67\% \\
        \bottomrule
    \end{tabular}}
\end{table}

Главный эффект дают групповые ограничения. Без них оптимизатор берет максимально разрешенное плечо и ставит почти все бумаги на верхнюю границу, кроме GAZP. После добавления групп портфель становится заметно спокойнее.

\begin{figure}[H]
    \centering
    \includegraphics[width=\textwidth]{efficient_frontiers.png}
    \caption{Эффективные границы для разных наборов ограничений}
\end{figure}

Эффективная граница с полным набором ограничений лежит ниже свободного варианта. Это нормально: чем больше ограничений, тем меньше допустимых портфелей и тем меньше максимум доходности при том же риске.

\subsection{Итоговый портфель}

Для основного варианта берем все ограничения. При $\theta=2,3,4$ получился один и тот же портфель.

\begin{table}[H]
    \centering
    \small
    \caption{Итоговый портфель при всех ограничениях}
    \resizebox{\textwidth}{!}{
    \begin{tabular}{lccccccccccc}
        \toprule
        $\theta$ & Режим & Доходность & Риск & $\xi$ & Безриск. доля & GAZP & ROSN & LKOH & FEES & SBER & VTBR \\
        \midrule
        2 & заем & 90.32\% & 20.53\% & 1.20 & -20.00\% & 5.00\% & 30.00\% & 6.33\% & 27.00\% & 39.00\% & 12.67\% \\
        3 & заем & 90.32\% & 20.53\% & 1.20 & -20.00\% & 5.00\% & 30.00\% & 6.33\% & 27.00\% & 39.00\% & 12.67\% \\
        4 & заем & 90.32\% & 20.53\% & 1.20 & -20.00\% & 5.00\% & 30.00\% & 6.33\% & 27.00\% & 39.00\% & 12.67\% \\
        \bottomrule
    \end{tabular}}
\end{table}

Здесь $\xi$ означает сумму долей в рискованных бумагах. Значение $\xi=1.20$ значит, что в акции вложено 120\% капитала, а еще 20\% профинансированы за счет займа.

\begin{figure}[H]
    \centering
    \includegraphics[width=\textwidth]{infographic_portfolio_weights.png}
    \caption{Состав оптимального портфеля при разных ограничениях}
\end{figure}

В итоговом портфеле виден понятный результат: GAZP стоит на минимальной границе, SBER стоит на максимальной, а остальные доли подстраиваются под групповые ограничения.

\subsection{Чувствительность к параметру риска}

Для проверки были посчитаны и более высокие значения $\theta$.

\begin{table}[H]
    \centering
    \caption{Как меняется портфель при росте неприятия риска}
    \begin{tabular}{ccccc}
        \toprule
        $\theta$ & Режим & Доходность & Риск & Безрисковая доля \\
        \midrule
        2 & заем & 90.32\% & 20.53\% & -20.00\% \\
        3 & заем & 90.32\% & 20.53\% & -20.00\% \\
        4 & заем & 90.32\% & 20.53\% & -20.00\% \\
        12 & заем & 89.80\% & 20.31\% & -20.00\% \\
        20 & заем & 81.50\% & 17.39\% & -5.93\% \\
        40 & вклад & 59.27\% & 12.58\% & 19.24\% \\
        \bottomrule
    \end{tabular}
\end{table}

\begin{figure}[H]
    \centering
    \includegraphics[width=\textwidth]{infographic_theta_sensitivity.png}
    \caption{Изменение риска, доходности и безрисковой доли при росте $\theta$}
\end{figure}

При $\theta=2,3,4$ портфель не меняется, потому что ограничения оказываются сильнее изменения параметра риска. При больших $\theta$ инвестор становится осторожнее: сначала снижается заем, а при $\theta=40$ часть капитала уже уходит в безрисковый актив.

\subsection{Активные ограничения}

Для итогового портфеля при $\theta=3$ активны следующие ограничения.

\begin{table}[H]
    \centering
    \caption{Активные ограничения при $\theta=3$}
    \begin{tabular}{llcc}
        \toprule
        Тип & Ограничение & Значение & Активная граница \\
        \midrule
        Акция & GAZP & 5.00\% & нижняя \\
        Акция & SBER & 39.00\% & верхняя \\
        Группа & Энергетика & 27.00\% & нижняя \\
        Группа & Внутренний рынок & 85.00\% & верхняя \\
        Группа & Внешний рынок & 35.00\% & верхняя \\
        Сравнение групп & Нефтегаз / банки & 0.80 & нижняя \\
        \bottomrule
    \end{tabular}
\end{table}

\begin{figure}[H]
    \centering
    \includegraphics[width=\textwidth]{infographic_active_constraints.png}
    \caption{Какие ограничения фактически формируют итоговый портфель}
\end{figure}

Эти ограничения и определяют форму решения. Если их ослабить, портфель снова начнет сильнее уходить в самые доходные бумаги.

\section{Выводы}

\begin{enumerate}
    \item По фактическому интервалу из файла самыми доходными по оценке оказались SBER и VTBR. У GAZP оценка доходности отрицательная.
    \item Без групповых ограничений оптимизатор берет слишком агрессивный портфель: почти все бумаги, кроме GAZP, ставятся на верхнюю границу.
    \item Групповые ограничения сильнее всего влияют на итоговый результат. Они ограничивают внешний рынок, внутренний рынок и минимальную долю энергетики.
    \item При $\theta=2,3,4$ оптимальный портфель одинаковый: 120\% капитала в рискованных активах и -20\% в безрисковом активе, то есть используется заем.
    \item При большом $\theta$ портфель становится осторожнее. При $\theta=40$ модель уже не занимает деньги, а держит около 19\% капитала в безрисковом активе.
\end{enumerate}

Итоговый портфель при основном расчете:
\[
\begin{array}{c|cccccc}
 & GAZP & ROSN & LKOH & FEES & SBER & VTBR \\
\hline
x_i & 5.00\% & 30.00\% & 6.33\% & 27.00\% & 39.00\% & 12.67\%
\end{array}
\]

\end{document}
